\documentclass[11pt,a4paper]{article}
\newcommand{\gradelabel}{Grade 8}
\newcommand{\leconhdr}{Déplacements simples}
% Préambule commun aux plans de cours de robotique (Grades 7 à 12).
% Chaque plan de cours commence par \documentclass[11pt,a4paper]{article}, définit
% \gradelabel et \leconhdr (en-têtes de page), puis % Préambule commun aux plans de cours de robotique (Grades 7 à 12).
% Chaque plan de cours commence par \documentclass[11pt,a4paper]{article}, définit
% \gradelabel et \leconhdr (en-têtes de page), puis % Préambule commun aux plans de cours de robotique (Grades 7 à 12).
% Chaque plan de cours commence par \documentclass[11pt,a4paper]{article}, définit
% \gradelabel et \leconhdr (en-têtes de page), puis \input{../../preambule_robotique.tex}.
\usepackage[T1]{fontenc}
\usepackage[utf8]{inputenc}
\usepackage{lmodern}
\usepackage[french]{babel}
\usepackage[margin=2cm,headheight=14pt]{geometry}
\usepackage[most]{tcolorbox}
\usepackage{xcolor,tabularx,booktabs,array,enumitem,fancyhdr,amssymb,tikz}
\usetikzlibrary{arrows.meta,positioning}

\pagestyle{fancy}
\fancyhf{}
\lhead{\small Robotique -- \gradelabel{} -- Novembre 2026}
\rhead{\small \leconhdr}
\cfoot{\thepage}
\setlist{itemsep=1pt,topsep=2pt,leftmargin=*}
\renewcommand{\arraystretch}{1.35}
\newcolumntype{Y}{>{\raggedright\arraybackslash}X}
\newcommand{\cb}{$\square$\,}
\newcommand{\fl}{$\rightarrow$}
\newcommand{\lig}{\noindent\rule{\linewidth}{0.3pt}\\[5mm]}

\newtcolorbox{etape}[3]{enhanced,breakable,colback=#1!6,colframe=#1!70!black,
  coltitle=white,fonttitle=\bfseries,title={#2\hfill #3},boxrule=0.8pt,arc=2mm,
  before skip=8pt,after skip=8pt}
\newtcolorbox{astuce}[1][Différenciation et besoins particuliers]{breakable,
  colback=yellow!12,colframe=yellow!50!black,fonttitle=\bfseries\small,
  title={#1},boxrule=0.5pt,arc=1mm,before skip=6pt,after skip=6pt}
\newtcolorbox{code}[1][Code]{breakable,colback=black!4,colframe=black!60,
  fonttitle=\bfseries\small,title={#1},boxrule=0.5pt,arc=1mm,fontupper=\ttfamily\small,
  before skip=6pt,after skip=6pt}
% Blocs de code coloré (écrire le code sans accents : comme dans l'éditeur)
\newtcblisting{arduino}[1][Code Arduino]{listing only,breakable,colback=black!4,
  colframe=black!60,fonttitle=\bfseries\small,title={#1},boxrule=0.5pt,arc=1mm,
  before skip=6pt,after skip=6pt,
  listing options={language=C++,basicstyle=\ttfamily\small,columns=fullflexible,
  keepspaces=true,showstringspaces=false,commentstyle=\color{green!40!black}\itshape,
  keywordstyle=\color{blue!70!black}\bfseries}}
\newtcblisting{python}[1][Code Python]{listing only,breakable,colback=black!4,
  colframe=black!60,fonttitle=\bfseries\small,title={#1},boxrule=0.5pt,arc=1mm,
  before skip=6pt,after skip=6pt,
  listing options={language=Python,basicstyle=\ttfamily\small,columns=fullflexible,
  keepspaces=true,showstringspaces=false,commentstyle=\color{green!40!black}\itshape,
  keywordstyle=\color{blue!70!black}\bfseries}}
\newcommand{\entete}[2]{%
  \begin{tcolorbox}[colback=#1!10,colframe=#1!70!black,boxrule=0.8pt,arc=2mm]
  {\Large\bfseries #2}\hfill Nom : \rule{4cm}{0.3pt}\quad Date : \rule{2.2cm}{0.3pt}
  \end{tcolorbox}}

% Bandeau de titre : #1 titre, #2 semaine
\newcommand{\titreblock}[2]{%
\begin{tcolorbox}[colback=teal!10,colframe=teal!70!black,arc=2mm,boxrule=1pt]
\centering
{\LARGE\bfseries #1}\\[2mm]
{\large Robotique -- \gradelabel{} -- Plan de cours de 2 heures}\\[1mm]
Enseignant : Francis Ki \quad|\quad #2
\end{tcolorbox}}

% Tableau du minutage : 6 contenus (accroche, mini-leçon, guidée, autonome, mise en commun, exit ticket)
\newcommand{\minutage}[6]{%
\begin{tabularx}{\linewidth}{>{\bfseries}p{2.2cm}p{2.6cm}Y p{3cm}}
\toprule
Minutes & Étape & Contenu & Organisation\\
\midrule
0 -- 10 & Accroche & #1 & Classe entière\\
10 -- 30 & Mini-leçon & #2 & Classe entière\\
30 -- 60 & Pratique guidée & #3 & Collectif puis binômes\\
60 -- 95 & Pratique autonome & #4 & Individuel sur poste\\
95 -- 110 & Mise en commun & #5 & Classe entière\\
110 -- 120 & Exit ticket & #6 & Individuel sur poste\\
\bottomrule
\end{tabularx}}

% Étayage progressif : #1 aide 1, #2 aide 2, #3 aide 3
\newcommand{\aides}[3]{%
\textbf{Étayage progressif de l'aide.}
\begin{center}
\begin{tabularx}{\linewidth}{>{\bfseries}p{2.4cm}Y p{4.4cm}}
\toprule
Niveau d'aide & Contenu & Qui intervient\\
\midrule
Aide 1 & #1 & Le binôme seul\\
Aide 2 & #2 & Le binôme avec la carte d'aide\\
Aide 3 & #3 & L'enseignant\\
\bottomrule
\end{tabularx}
\end{center}}

% Les trois groupes de la pratique autonome : #1 tâche A, #2 tâche B, #3 tâche C
\newcommand{\groupes}[3]{%
\textbf{Consigne commune\,:} chaque élève ouvre le fichier de son groupe, le renomme avec son prénom et travaille seul sur son poste. L'enseignant passe d'abord auprès du groupe A pendant 10 min.

\medskip
\begin{tabularx}{\linewidth}{>{\bfseries}p{3cm}Y p{2.2cm}}
\toprule
Groupe & Tâche sur poste & Annexe\\
\midrule
A : en dessous du niveau & #1 & Annexe C\\
B : au niveau & #2 & Annexe D\\
C : au-dessus du niveau & #3 & Annexe E\\
\bottomrule
\end{tabularx}}

% Mise en commun : #1 groupe A, #2 groupe B, #3 groupe C, #4 conclusion
\newcommand{\miseencommun}[4]{%
\begin{etape}{teal}{5. Mise en commun}{95 -- 110 min}
\begin{enumerate}
 \item (2 min) Retour au calme, écrans mis en veille.
 \item (9 min) Trois passages courts de 3 min, avec valorisation :
 \begin{itemize}
  \item \textbf{Groupe A} : #1
  \item \textbf{Groupe B} : #2
  \item \textbf{Groupe C} : #3
 \end{itemize}
 \item (4 min) Conclusion : #4
\end{enumerate}
\end{etape}}

% Clôture : #1 question 1, #2 question 2
\newcommand{\cloture}[2]{%
\begin{etape}{gray}{6. Clôture : exit ticket}{110 -- 120 min}
Formulaire en ligne ou papier (annexe G) :
\begin{enumerate}
 \item #1
 \item #2
 \item Auto-évaluation : \emph{Je maîtrise / Je progresse / J'ai besoin d'aide.}
\end{enumerate}
Lire les réponses avant la séance suivante pour ajuster les groupes de besoin.
\end{etape}}

% Grille d'observation : #1 #2 #3 intitulés de colonnes
\newcommand{\grille}[3]{%
\subsection*{Grille d'observation (enseignant)}
\begin{tabularx}{\linewidth}{Y p{2.4cm} p{2.4cm} p{2.4cm} Y}
\toprule
Élève & #1 & #2 & #3 & Groupe suivant\\
\midrule
 & & & & \\ & & & & \\ & & & & \\ & & & & \\ & & & & \\
\bottomrule
\end{tabularx}}

% Signature de l'exit ticket
\newcommand{\autoeval}{Je me sens : \cb Je maîtrise \quad \cb Je progresse \quad \cb J'ai besoin d'aide}
.
\usepackage[T1]{fontenc}
\usepackage[utf8]{inputenc}
\usepackage{lmodern}
\usepackage[french]{babel}
\usepackage[margin=2cm,headheight=14pt]{geometry}
\usepackage[most]{tcolorbox}
\usepackage{xcolor,tabularx,booktabs,array,enumitem,fancyhdr,amssymb,tikz}
\usetikzlibrary{arrows.meta,positioning}

\pagestyle{fancy}
\fancyhf{}
\lhead{\small Robotique -- \gradelabel{} -- Novembre 2026}
\rhead{\small \leconhdr}
\cfoot{\thepage}
\setlist{itemsep=1pt,topsep=2pt,leftmargin=*}
\renewcommand{\arraystretch}{1.35}
\newcolumntype{Y}{>{\raggedright\arraybackslash}X}
\newcommand{\cb}{$\square$\,}
\newcommand{\fl}{$\rightarrow$}
\newcommand{\lig}{\noindent\rule{\linewidth}{0.3pt}\\[5mm]}

\newtcolorbox{etape}[3]{enhanced,breakable,colback=#1!6,colframe=#1!70!black,
  coltitle=white,fonttitle=\bfseries,title={#2\hfill #3},boxrule=0.8pt,arc=2mm,
  before skip=8pt,after skip=8pt}
\newtcolorbox{astuce}[1][Différenciation et besoins particuliers]{breakable,
  colback=yellow!12,colframe=yellow!50!black,fonttitle=\bfseries\small,
  title={#1},boxrule=0.5pt,arc=1mm,before skip=6pt,after skip=6pt}
\newtcolorbox{code}[1][Code]{breakable,colback=black!4,colframe=black!60,
  fonttitle=\bfseries\small,title={#1},boxrule=0.5pt,arc=1mm,fontupper=\ttfamily\small,
  before skip=6pt,after skip=6pt}
% Blocs de code coloré (écrire le code sans accents : comme dans l'éditeur)
\newtcblisting{arduino}[1][Code Arduino]{listing only,breakable,colback=black!4,
  colframe=black!60,fonttitle=\bfseries\small,title={#1},boxrule=0.5pt,arc=1mm,
  before skip=6pt,after skip=6pt,
  listing options={language=C++,basicstyle=\ttfamily\small,columns=fullflexible,
  keepspaces=true,showstringspaces=false,commentstyle=\color{green!40!black}\itshape,
  keywordstyle=\color{blue!70!black}\bfseries}}
\newtcblisting{python}[1][Code Python]{listing only,breakable,colback=black!4,
  colframe=black!60,fonttitle=\bfseries\small,title={#1},boxrule=0.5pt,arc=1mm,
  before skip=6pt,after skip=6pt,
  listing options={language=Python,basicstyle=\ttfamily\small,columns=fullflexible,
  keepspaces=true,showstringspaces=false,commentstyle=\color{green!40!black}\itshape,
  keywordstyle=\color{blue!70!black}\bfseries}}
\newcommand{\entete}[2]{%
  \begin{tcolorbox}[colback=#1!10,colframe=#1!70!black,boxrule=0.8pt,arc=2mm]
  {\Large\bfseries #2}\hfill Nom : \rule{4cm}{0.3pt}\quad Date : \rule{2.2cm}{0.3pt}
  \end{tcolorbox}}

% Bandeau de titre : #1 titre, #2 semaine
\newcommand{\titreblock}[2]{%
\begin{tcolorbox}[colback=teal!10,colframe=teal!70!black,arc=2mm,boxrule=1pt]
\centering
{\LARGE\bfseries #1}\\[2mm]
{\large Robotique -- \gradelabel{} -- Plan de cours de 2 heures}\\[1mm]
Enseignant : Francis Ki \quad|\quad #2
\end{tcolorbox}}

% Tableau du minutage : 6 contenus (accroche, mini-leçon, guidée, autonome, mise en commun, exit ticket)
\newcommand{\minutage}[6]{%
\begin{tabularx}{\linewidth}{>{\bfseries}p{2.2cm}p{2.6cm}Y p{3cm}}
\toprule
Minutes & Étape & Contenu & Organisation\\
\midrule
0 -- 10 & Accroche & #1 & Classe entière\\
10 -- 30 & Mini-leçon & #2 & Classe entière\\
30 -- 60 & Pratique guidée & #3 & Collectif puis binômes\\
60 -- 95 & Pratique autonome & #4 & Individuel sur poste\\
95 -- 110 & Mise en commun & #5 & Classe entière\\
110 -- 120 & Exit ticket & #6 & Individuel sur poste\\
\bottomrule
\end{tabularx}}

% Étayage progressif : #1 aide 1, #2 aide 2, #3 aide 3
\newcommand{\aides}[3]{%
\textbf{Étayage progressif de l'aide.}
\begin{center}
\begin{tabularx}{\linewidth}{>{\bfseries}p{2.4cm}Y p{4.4cm}}
\toprule
Niveau d'aide & Contenu & Qui intervient\\
\midrule
Aide 1 & #1 & Le binôme seul\\
Aide 2 & #2 & Le binôme avec la carte d'aide\\
Aide 3 & #3 & L'enseignant\\
\bottomrule
\end{tabularx}
\end{center}}

% Les trois groupes de la pratique autonome : #1 tâche A, #2 tâche B, #3 tâche C
\newcommand{\groupes}[3]{%
\textbf{Consigne commune\,:} chaque élève ouvre le fichier de son groupe, le renomme avec son prénom et travaille seul sur son poste. L'enseignant passe d'abord auprès du groupe A pendant 10 min.

\medskip
\begin{tabularx}{\linewidth}{>{\bfseries}p{3cm}Y p{2.2cm}}
\toprule
Groupe & Tâche sur poste & Annexe\\
\midrule
A : en dessous du niveau & #1 & Annexe C\\
B : au niveau & #2 & Annexe D\\
C : au-dessus du niveau & #3 & Annexe E\\
\bottomrule
\end{tabularx}}

% Mise en commun : #1 groupe A, #2 groupe B, #3 groupe C, #4 conclusion
\newcommand{\miseencommun}[4]{%
\begin{etape}{teal}{5. Mise en commun}{95 -- 110 min}
\begin{enumerate}
 \item (2 min) Retour au calme, écrans mis en veille.
 \item (9 min) Trois passages courts de 3 min, avec valorisation :
 \begin{itemize}
  \item \textbf{Groupe A} : #1
  \item \textbf{Groupe B} : #2
  \item \textbf{Groupe C} : #3
 \end{itemize}
 \item (4 min) Conclusion : #4
\end{enumerate}
\end{etape}}

% Clôture : #1 question 1, #2 question 2
\newcommand{\cloture}[2]{%
\begin{etape}{gray}{6. Clôture : exit ticket}{110 -- 120 min}
Formulaire en ligne ou papier (annexe G) :
\begin{enumerate}
 \item #1
 \item #2
 \item Auto-évaluation : \emph{Je maîtrise / Je progresse / J'ai besoin d'aide.}
\end{enumerate}
Lire les réponses avant la séance suivante pour ajuster les groupes de besoin.
\end{etape}}

% Grille d'observation : #1 #2 #3 intitulés de colonnes
\newcommand{\grille}[3]{%
\subsection*{Grille d'observation (enseignant)}
\begin{tabularx}{\linewidth}{Y p{2.4cm} p{2.4cm} p{2.4cm} Y}
\toprule
Élève & #1 & #2 & #3 & Groupe suivant\\
\midrule
 & & & & \\ & & & & \\ & & & & \\ & & & & \\ & & & & \\
\bottomrule
\end{tabularx}}

% Signature de l'exit ticket
\newcommand{\autoeval}{Je me sens : \cb Je maîtrise \quad \cb Je progresse \quad \cb J'ai besoin d'aide}
.
\usepackage[T1]{fontenc}
\usepackage[utf8]{inputenc}
\usepackage{lmodern}
\usepackage[french]{babel}
\usepackage[margin=2cm,headheight=14pt]{geometry}
\usepackage[most]{tcolorbox}
\usepackage{xcolor,tabularx,booktabs,array,enumitem,fancyhdr,amssymb,tikz}
\usetikzlibrary{arrows.meta,positioning}

\pagestyle{fancy}
\fancyhf{}
\lhead{\small Robotique -- \gradelabel{} -- Novembre 2026}
\rhead{\small \leconhdr}
\cfoot{\thepage}
\setlist{itemsep=1pt,topsep=2pt,leftmargin=*}
\renewcommand{\arraystretch}{1.35}
\newcolumntype{Y}{>{\raggedright\arraybackslash}X}
\newcommand{\cb}{$\square$\,}
\newcommand{\fl}{$\rightarrow$}
\newcommand{\lig}{\noindent\rule{\linewidth}{0.3pt}\\[5mm]}

\newtcolorbox{etape}[3]{enhanced,breakable,colback=#1!6,colframe=#1!70!black,
  coltitle=white,fonttitle=\bfseries,title={#2\hfill #3},boxrule=0.8pt,arc=2mm,
  before skip=8pt,after skip=8pt}
\newtcolorbox{astuce}[1][Différenciation et besoins particuliers]{breakable,
  colback=yellow!12,colframe=yellow!50!black,fonttitle=\bfseries\small,
  title={#1},boxrule=0.5pt,arc=1mm,before skip=6pt,after skip=6pt}
\newtcolorbox{code}[1][Code]{breakable,colback=black!4,colframe=black!60,
  fonttitle=\bfseries\small,title={#1},boxrule=0.5pt,arc=1mm,fontupper=\ttfamily\small,
  before skip=6pt,after skip=6pt}
% Blocs de code coloré (écrire le code sans accents : comme dans l'éditeur)
\newtcblisting{arduino}[1][Code Arduino]{listing only,breakable,colback=black!4,
  colframe=black!60,fonttitle=\bfseries\small,title={#1},boxrule=0.5pt,arc=1mm,
  before skip=6pt,after skip=6pt,
  listing options={language=C++,basicstyle=\ttfamily\small,columns=fullflexible,
  keepspaces=true,showstringspaces=false,commentstyle=\color{green!40!black}\itshape,
  keywordstyle=\color{blue!70!black}\bfseries}}
\newtcblisting{python}[1][Code Python]{listing only,breakable,colback=black!4,
  colframe=black!60,fonttitle=\bfseries\small,title={#1},boxrule=0.5pt,arc=1mm,
  before skip=6pt,after skip=6pt,
  listing options={language=Python,basicstyle=\ttfamily\small,columns=fullflexible,
  keepspaces=true,showstringspaces=false,commentstyle=\color{green!40!black}\itshape,
  keywordstyle=\color{blue!70!black}\bfseries}}
\newcommand{\entete}[2]{%
  \begin{tcolorbox}[colback=#1!10,colframe=#1!70!black,boxrule=0.8pt,arc=2mm]
  {\Large\bfseries #2}\hfill Nom : \rule{4cm}{0.3pt}\quad Date : \rule{2.2cm}{0.3pt}
  \end{tcolorbox}}

% Bandeau de titre : #1 titre, #2 semaine
\newcommand{\titreblock}[2]{%
\begin{tcolorbox}[colback=teal!10,colframe=teal!70!black,arc=2mm,boxrule=1pt]
\centering
{\LARGE\bfseries #1}\\[2mm]
{\large Robotique -- \gradelabel{} -- Plan de cours de 2 heures}\\[1mm]
Enseignant : Francis Ki \quad|\quad #2
\end{tcolorbox}}

% Tableau du minutage : 6 contenus (accroche, mini-leçon, guidée, autonome, mise en commun, exit ticket)
\newcommand{\minutage}[6]{%
\begin{tabularx}{\linewidth}{>{\bfseries}p{2.2cm}p{2.6cm}Y p{3cm}}
\toprule
Minutes & Étape & Contenu & Organisation\\
\midrule
0 -- 10 & Accroche & #1 & Classe entière\\
10 -- 30 & Mini-leçon & #2 & Classe entière\\
30 -- 60 & Pratique guidée & #3 & Collectif puis binômes\\
60 -- 95 & Pratique autonome & #4 & Individuel sur poste\\
95 -- 110 & Mise en commun & #5 & Classe entière\\
110 -- 120 & Exit ticket & #6 & Individuel sur poste\\
\bottomrule
\end{tabularx}}

% Étayage progressif : #1 aide 1, #2 aide 2, #3 aide 3
\newcommand{\aides}[3]{%
\textbf{Étayage progressif de l'aide.}
\begin{center}
\begin{tabularx}{\linewidth}{>{\bfseries}p{2.4cm}Y p{4.4cm}}
\toprule
Niveau d'aide & Contenu & Qui intervient\\
\midrule
Aide 1 & #1 & Le binôme seul\\
Aide 2 & #2 & Le binôme avec la carte d'aide\\
Aide 3 & #3 & L'enseignant\\
\bottomrule
\end{tabularx}
\end{center}}

% Les trois groupes de la pratique autonome : #1 tâche A, #2 tâche B, #3 tâche C
\newcommand{\groupes}[3]{%
\textbf{Consigne commune\,:} chaque élève ouvre le fichier de son groupe, le renomme avec son prénom et travaille seul sur son poste. L'enseignant passe d'abord auprès du groupe A pendant 10 min.

\medskip
\begin{tabularx}{\linewidth}{>{\bfseries}p{3cm}Y p{2.2cm}}
\toprule
Groupe & Tâche sur poste & Annexe\\
\midrule
A : en dessous du niveau & #1 & Annexe C\\
B : au niveau & #2 & Annexe D\\
C : au-dessus du niveau & #3 & Annexe E\\
\bottomrule
\end{tabularx}}

% Mise en commun : #1 groupe A, #2 groupe B, #3 groupe C, #4 conclusion
\newcommand{\miseencommun}[4]{%
\begin{etape}{teal}{5. Mise en commun}{95 -- 110 min}
\begin{enumerate}
 \item (2 min) Retour au calme, écrans mis en veille.
 \item (9 min) Trois passages courts de 3 min, avec valorisation :
 \begin{itemize}
  \item \textbf{Groupe A} : #1
  \item \textbf{Groupe B} : #2
  \item \textbf{Groupe C} : #3
 \end{itemize}
 \item (4 min) Conclusion : #4
\end{enumerate}
\end{etape}}

% Clôture : #1 question 1, #2 question 2
\newcommand{\cloture}[2]{%
\begin{etape}{gray}{6. Clôture : exit ticket}{110 -- 120 min}
Formulaire en ligne ou papier (annexe G) :
\begin{enumerate}
 \item #1
 \item #2
 \item Auto-évaluation : \emph{Je maîtrise / Je progresse / J'ai besoin d'aide.}
\end{enumerate}
Lire les réponses avant la séance suivante pour ajuster les groupes de besoin.
\end{etape}}

% Grille d'observation : #1 #2 #3 intitulés de colonnes
\newcommand{\grille}[3]{%
\subsection*{Grille d'observation (enseignant)}
\begin{tabularx}{\linewidth}{Y p{2.4cm} p{2.4cm} p{2.4cm} Y}
\toprule
Élève & #1 & #2 & #3 & Groupe suivant\\
\midrule
 & & & & \\ & & & & \\ & & & & \\ & & & & \\ & & & & \\
\bottomrule
\end{tabularx}}

% Signature de l'exit ticket
\newcommand{\autoeval}{Je me sens : \cb Je maîtrise \quad \cb Je progresse \quad \cb J'ai besoin d'aide}


\begin{document}
\titreblock{Déplacements simples}{Semaine du 9 au 13 novembre 2026}

%==================================================================
\section{Informations générales}
\begin{tabularx}{\linewidth}{>{\bfseries}p{3.6cm}Y}
\toprule
Titre & Déplacements simples du mBot\\
Classe / effectif & Grade 8, 11 à 20 élèves, niveaux hétérogènes, élèves à besoins particuliers\\
Durée & 2 heures (120 min)\\
Lieu & Salle informatique, un poste par élève\\
Matériel et logiciels &
\begin{itemize}
 \item Un poste par élève avec mBlock 5 (connexion robot opérationnelle depuis la séance précédente).
 \item Un robot mBot par binôme, chargé.
 \item Un mètre ruban ou une règle de 1 m par binôme, du ruban de masquage pour marquer le départ.
 \item Un tableur (Google Sheets ou LibreOffice Calc) pour le niveau C.
\end{itemize}\\
Pré-requis &
Séance 1 : connexion du mBot, LED, buzzer, bouton. Savoir mesurer une longueur en centimètres.\\
\bottomrule
\end{tabularx}

\begin{astuce}[Préparation avant le cours (environ 40 min)]
\begin{enumerate}
 \item Tracer une piste de 2 m au sol avec une graduation tous les 10 cm (ruban).
 \item Tester les robots : certains peuvent dévier légèrement vers un côté. Noter les robots à problème.
 \item Préparer les fichiers de départ (annexes C à E) et le formulaire de sortie.
 \item Vérifier le niveau des batteries : une batterie faible fausse les mesures.
\end{enumerate}
\end{astuce}

\begin{astuce}[Sécurité]
Les mesures se font au sol, pas sur la table. Pas d'élève devant un robot en marche. Toujours un bloc \texttt{arrêter} en fin de programme.
\end{astuce}

%==================================================================
\section{Objectifs d'apprentissage}
\begin{tcolorbox}[colback=gray!8,colframe=gray!60!black,arc=1mm,boxrule=0.5pt]
\textbf{Objectifs du programme}
\begin{itemize}
 \item Programmer un déplacement rectiligne.
 \item Contrôler la vitesse et la durée.
\end{itemize}
\textbf{À la fin de la séance, l'élève est capable\,:}
\begin{itemize}
 \item d'écrire un programme \texttt{avancer, attendre, arrêter}\,;
 \item d'expliquer l'effet de la puissance et de la durée sur la distance parcourue\,;
 \item de calibrer le robot par essais pour atteindre une distance donnée.
\end{itemize}
\textbf{Compétences numériques\,:} mesurer, relever des valeurs dans un tableau, ajuster un programme par essais.
\end{tcolorbox}

%==================================================================
\section{Déroulé minuté (120 min)}
\minutage{Course de robots humains}{Puissance, durée et distance}{Un essai mesuré en collectif, puis le tableau puissance et durée en binômes}{Trois niveaux de besoin}{Valoriser chaque groupe}{Trois questions sur poste}

%------------------------------------------------------------------
\begin{etape}{orange}{1. Accroche : course de robots humains}{0 -- 10 min}
\begin{enumerate}
 \item (6 min) Trois élèves jouent des robots : l'un marche lentement, l'autre vite, le troisième vite mais seulement 2 secondes. Chronométrer et mesurer la distance parcourue. Qui va le plus loin ?
 \item (4 min) Question-problème : \textbf{«~Comment faire avancer un robot exactement de 1 mètre ?~»} Noter les idées (vitesse, temps).
\end{enumerate}
\end{etape}

%------------------------------------------------------------------
\begin{etape}{blue}{2. Mini-leçon : puissance, durée, distance}{10 -- 30 min}
\textbf{Support\,:} annexe A. Quatre temps de 5 min.

\textbf{Temps 1 : le programme de base.}
\begin{code}[Programme de déplacement]
avancer à la puissance 50 \%\\
attendre 2 secondes\\
arrêter les moteurs
\end{code}
Sans le bloc \emph{arrêter}, le robot continue d'avancer. Les noms exacts des blocs dépendent de la version de mBlock.

\textbf{Temps 2 : deux variables de réglage.} La \textbf{puissance} (en pourcentage) règle la vitesse. La \textbf{durée} règle le temps de marche.

\textbf{Temps 3 : la relation.} Plus la puissance et la durée sont grandes, plus le robot va loin. La relation n'est pas parfaitement proportionnelle (démarrage, batterie, sol). Il faut donc \emph{mesurer}.

\begin{center}
\begin{tikzpicture}[box/.style={draw,rounded corners,thick,minimum height=1cm,text width=2.8cm,align=center,font=\small},>={Stealth[length=3mm]},node distance=0.8cm]
\node[box,fill=orange!20] (a) {Puissance};
\node[box,fill=blue!15,right=of a] (b) {Durée};
\node[box,fill=green!15,right=of b] (c) {Distance parcourue};
\draw[->,thick] (a)--(c); \draw[->,thick] (b)--(c);
\end{tikzpicture}
\end{center}

\textbf{Temps 4 : la démarche d'essais.} Mesurer, comparer à la cible, corriger la durée, recommencer. Écrire chaque essai dans un tableau.

\begin{astuce}
Fournir un tableau d'essais pré-imprimé. Autoriser la calculatrice. Faire répéter la démarche : mesurer, comparer, corriger.
\end{astuce}
\end{etape}

%------------------------------------------------------------------
\begin{etape}{green}{3. Pratique guidée}{30 -- 60 min}
\textbf{Fiche\,:} annexe B.

\textbf{3a. Un essai mesuré en collectif (12 min).} Un robot (celui de l'enseignant, projeté par caméra ou en démonstration) : puissance 50 \%, durée 2 s. Mesurer la distance sur la piste. Chaque binôme note la valeur.

\textbf{3b. Le tableau puissance et durée en binômes (18 min).} Chaque binôme réalise six essais (puissances 30, 50 et 70 \%, durées 1 et 2 s) et remplit le tableau B2. Alternance des rôles : un pilote, un mesureur.

\aides{Le programme du 3a et le tableau d'essais à moitié rempli}{Questions guides : \emph{Que change-t-on ? Qu'est-ce qui reste identique ? Que se passe-t-il quand la durée double ?}}{L'enseignant réalise un essai avec le binôme, qui continue}

\begin{astuce}
Attribuer un rôle précis à chaque élève du binôme (pilote, mesureur, enregistreur) et les faire tourner. Accepter des mesures approchées à 2 cm près.
\end{astuce}
\end{etape}

%------------------------------------------------------------------
\begin{etape}{violet}{4. Pratique autonome différenciée}{60 -- 95 min}
Chaque élève prépare son programme seul sur son poste. Les essais se font par binôme sur la piste, à tour de rôle.

\medskip
\groupes{Programme donné : avancer à 50 \% pendant 2 s puis arrêter. Mesurer la distance, puis changer \emph{uniquement} la durée pour atteindre environ 50 cm, avec une banque de durées (1, 1,5, 2, 3).}{Faire avancer le robot de 1 mètre (à 10 cm près) en choisissant puissance et durée par essais. Tenir un tableau d'essais.}{Programmer un aller-retour : avancer de 80 cm, s'arrêter, reculer pour revenir au départ à 5 cm près. Reporter les mesures dans un tableur et tracer la distance en fonction de la durée.}

\medskip
\textbf{Plan B si un robot dévie\,:} corriger en réglant légèrement les puissances des deux moteurs (par exemple 52 \% et 50 \%) ou changer de robot.
\end{etape}

%------------------------------------------------------------------
\miseencommun{un élève montre son tableau d'essais et lit son meilleur résultat. Valoriser la démarche plus que le résultat.}{un binôme montre comment il a atteint 1 mètre en deux ou trois essais. Valoriser la correction méthodique.}{un élève présente son graphique et explique ce qu'il remarque. Valoriser l'analyse.}{\textit{pour atteindre une distance, on règle la puissance et la durée et on corrige par essais.}}

\cloture{Quel bloc faut-il toujours mettre à la fin d'un déplacement ?}{Si ton robot va trop près de la cible, que modifies-tu ?}

%==================================================================
% ANNEXES
%==================================================================
\newpage
\section*{Annexe A : Support résumé de la mini-leçon}
\entete{blue}{Déplacements simples}

\begin{tcolorbox}[colback=yellow!10,colframe=yellow!50!black]
\textbf{Distance parcourue} dépend de la \textbf{puissance} (vitesse) et de la \textbf{durée}. Il faut \textbf{mesurer} puis \textbf{corriger}.
\end{tcolorbox}

\begin{code}[Programme de base]
avancer à la puissance 50 \%\\
attendre 2 secondes\\
arrêter les moteurs
\end{code}

\begin{tabularx}{\linewidth}{>{\bfseries}p{3.2cm}Y}
\toprule
Étape & Ce que je fais\\
\midrule
1 & Je prévois une valeur\\
2 & Je mesure la distance\\
3 & Je compare à la cible\\
4 & Je corrige la durée ou la puissance\\
\bottomrule
\end{tabularx}

%==================================================================
\newpage
\section*{Annexe B : Fiche de pratique guidée}
\entete{green}{Pratique guidée : puissance, durée et distance}

\subsection*{3a. Essai collectif}
Puissance : \dotfill\quad Durée : \dotfill\quad Distance mesurée : \dotfill

\subsection*{3b. Tableau d'essais (en binôme)}
\begin{tabularx}{\linewidth}{>{\bfseries}p{3cm}Y Y Y}
\toprule
Essai & Puissance (\%) & Durée (s) & Distance (cm)\\
\midrule
1 & 30 & 1 & \\
2 & 30 & 2 & \\
3 & 50 & 1 & \\
4 & 50 & 2 & \\
5 & 70 & 1 & \\
6 & 70 & 2 & \\
\bottomrule
\end{tabularx}

Ce que je remarque : \dotfill

\begin{tcolorbox}[colback=yellow!10,colframe=yellow!50!black,title={Cartes d'aide à utiliser dans l'ordre},fonttitle=\bfseries\small]
\textbf{Aide 1 :} regarde le programme et le tableau à moitié rempli.\\
\textbf{Aide 2 :} Que change-t-on ? Qu'est-ce qui reste identique ? Que se passe-t-il quand la durée double ?\\
\textbf{Aide 3 :} demande à l'enseignant de faire un essai avec toi.
\end{tcolorbox}

%==================================================================
\newpage
\section*{Annexe C : Niveau A (en dessous du niveau)}
\entete{orange}{Niveau A : j'ajuste la durée}

\textbf{Banque de durées :} 1 s, 1,5 s, 2 s, 3 s.

\subsection*{Activité 1 : essai de départ}
Puissance 50 \%, durée 2 s. Distance : \dotfill{} cm.

\subsection*{Activité 2 : cible de 50 cm}
\begin{tabularx}{\linewidth}{>{\bfseries}p{2.4cm}Y Y Y}
\toprule
Essai & Durée choisie & Distance mesurée & Trop court / trop long ?\\
\midrule
1 & & & \\
2 & & & \\
3 & & & \\
\bottomrule
\end{tabularx}

\subsection*{Activité 3 : complète}
Si le robot va trop loin, je \dotfill{} la durée. S'il ne va pas assez loin, je \dotfill{} la durée.

%==================================================================
\newpage
\section*{Annexe D : Niveau B (au niveau)}
\entete{blue}{Niveau B : atteindre 1 mètre}

\textbf{Cible :} 100 cm à 10 cm près.

\begin{tabularx}{\linewidth}{>{\bfseries}p{2cm}Y Y Y Y}
\toprule
Essai & Puissance (\%) & Durée (s) & Distance (cm) & Écart à la cible\\
\midrule
1 & & & & \\ 2 & & & & \\ 3 & & & & \\ 4 & & & & \\ 5 & & & & \\
\bottomrule
\end{tabularx}

\subsection*{Questions}
\begin{enumerate}
 \item Combien d'essais as-tu eu besoin ? \dotfill
 \item Pourquoi deux robots avec le même programme ne vont-ils pas exactement aussi loin ? \dotfill
\end{enumerate}

%==================================================================
\newpage
\section*{Annexe E : Niveau C (au-dessus du niveau)}
\entete{violet}{Niveau C : aller-retour précis}

\textbf{Mission :} avancer de 80 cm, s'arrêter, reculer et revenir au point de départ à 5 cm près.

\begin{tabularx}{\linewidth}{>{\bfseries}p{2cm}Y Y Y Y}
\toprule
Essai & Durée avant (s) & Durée arrière (s) & Écart au départ (cm) & Remarque\\
\midrule
1 & & & & \\ 2 & & & & \\ 3 & & & & \\
\bottomrule
\end{tabularx}

\subsection*{Graphique}
Dans le tableur, saisis la durée (en abscisse) et la distance (en ordonnée) pour une même puissance, puis trace le nuage de points.\\
Que remarques-tu ? \lig

\textbf{Défi bonus :} compare deux puissances sur le même graphique.

%==================================================================
\newpage
\section*{Annexe F : Corrections et attendus (enseignant)}

\subsection*{Annexe B}
Plus la puissance et la durée augmentent, plus la distance augmente. Quand la durée double, la distance double à peu près (pas exactement). Les valeurs dépendent du robot et du sol.

\subsection*{Annexe C}
Activité 3 : \textbf{diminuer} puis \textbf{augmenter}.

\subsection*{Annexe D}
Méthode attendue : essai, mesure, écart, correction proportionnelle (si 50 cm pour 2 s, alors environ 4 s pour 100 cm). Les robots diffèrent à cause des moteurs, du sol, de la batterie.

\subsection*{Annexe E}
Le nuage de points est à peu près aligné : la distance augmente avec la durée. L'écart de retour vient des frottements et d'un démarrage non instantané.\\
\textbf{Critères de réussite :} tableau d'essais complet, correction argumentée, graphique correct.

%==================================================================
\newpage
\section*{Annexe G : Exit ticket et grille d'observation}
\entete{gray}{Exit ticket : déplacements simples}

\begin{enumerate}
 \item Quel bloc faut-il toujours mettre à la fin d'un déplacement ?\\ \lig
 \item Si ton robot va trop près de la cible, que modifies-tu ?\\ \lig
 \item \autoeval
\end{enumerate}
\textbf{Attendus :} 1. \texttt{arrêter les moteurs}. 2. Augmenter la durée (ou la puissance).

\bigskip
\grille{Programme complet}{Tableau d'essais}{Correction argumentée}

\end{document}
